\chapter{Model Implementation}
\label{chap:implementation}

% ============================================================================
% CHAPTER PURPOSE
% Explain how the Chapter \ref{chap:model} formulation was implemented as a
% reproducible analysis tool. Software engineering is relevant where it protects
% calculation correctness, traceability, experimentation or usability; avoid a
% generic catalogue of frameworks and coding practices.
% ============================================================================

\section{Implementation Requirements}
\label{sec:implementation-requirements}

% HIGH PRIORITY. Translate research objectives into functional requirements:
% - load/select historic meeting, session and drivers;
% - obtain/prepare lap, stint, compound and pit information;
% - configure assumptions/parameters;
% - execute both scenarios for candidate decision laps;
% - display numerical and race trace outputs;
% - retain enough metadata for reproduction.
% Add non-functional requirements only where relevant: determinism, separation
% of model/UI, responsiveness, offline analysis and clear error handling.

\section{Technology and Architecture Selection}
\label{sec:technology-selection}

\subsection{Initial Web-Based Approach}
% Briefly describe the original Blazor/client-server proposal and what research
% need it was intended to satisfy. Do not teach basic client-server architecture.
% The generic web architecture figure should be removed unless directly useful.

\subsection{Transition to WPF}
% Preserve the existing rationale but tighten it into an engineering decision:
% - limited project time and front-end experience;
% - unreliable debugging/development setup;
% - WPF enabled focus on model functionality and repeatable desktop analysis;
% - trade-off: Windows-only tool and reduced deployment flexibility.
% Avoid self-deprecating wording. Demonstrate risk management and scope control.

\subsection{Final Architecture}
% SUGGESTED FIGURE 4.1 -- MUST HAVE:
% Data provider -> preprocessing/domain data -> model engine -> result objects ->
% race trace/UI. Show where user parameters enter.
% Explain separation of model engine from presentation so equations can be unit
% tested and reused without the WPF UI.

\section{Historic Data Acquisition}
\label{sec:data-acquisition}

\subsection{Data-Source Evaluation}
% Clarify the final role of FastF1 and OpenF1. Earlier drafts mention moving from
% FastF1 to OpenF1, while appendices still list FastF1. Resolve this explicitly:
% - what each source was evaluated for;
% - which is used in the final implementation;
% - whether both are used for different responsibilities;
% - licensing, rate limits, availability and fields required.
% Cite official documentation/repositories and record exact versions/endpoints.

\subsection{Meeting, Session and Driver Selection}
% Explain mapping from championship/event/session to driver records.
% Include sprint/race session handling if implemented.
% LOW/MEDIUM PRIORITY screenshots can illustrate workflow, but avoid turning
% this into a user manual.

\subsection{Lap, Stint and Pit Stop Preparation}
% HIGH PRIORITY. Describe transformations that affect model validity:
% - invalid/incomplete laps;
% - pit in/out laps;
% - Safety Car/VSC filtering;
% - compound and tyre-age reconstruction;
% - missing sectors/lap times;
% - data ordering and synchronisation;
% - whether laps are fuel-corrected or left as observed.
% This is engineering methodology, not merely implementation detail.

\subsection{Caching and Reproducibility}
% MEDIUM PRIORITY. If API responses are cached/localised, explain how this
% protects reproducibility and avoids data changes/rate limits.
% Record session identifiers and parameter sets used in validation.

\section{Model Engine Implementation}
\label{sec:model-engine-implementation}

\subsection{Domain Representation}
% Explain the principal model entities conceptually: driver state, tyre state,
% scenario, lap prediction and comparison result. A small class/domain diagram
% is useful only if it clarifies correspondence with Chapter 3 notation.

\subsection{Scenario Calculation Workflow}
% Map implementation steps directly to equations/sections:
% 1. build decision state;
% 2. predict each scenario lap;
% 3. apply optional penalties;
% 4. accumulate to common endpoint;
% 5. generate gain/post stop metrics;
% 6. repeat across candidate laps/offsets.
% Include pseudocode or a flowchart rather than large C# listings.

\subsection{Parameter Configuration and Traceability}
% Explain how values are entered/derived and associated with a run.
% Show validation of units/ranges and handling of missing values.
% Consider a parameter-summary panel/export if implemented.

\subsection{Error Handling}
% MEDIUM PRIORITY. Cover only failures that affect engineering confidence:
% unavailable API, incomplete session, invalid tyre age, no comparison window,
% or impossible parameter values. State whether the user is warned or the case
% is excluded; never silently substitute unsupported data.

\section{Race Trace Visualisation}
\label{sec:race-trace}

\subsection{Race Trace Definition}
% Briefly explain axes/reference convention for a reader unfamiliar with race
% traces. State exactly how gap is calculated and what the reference line means.

\subsection{Opportunity and Event Annotations}
% VERY HIGH VALUE to practical contribution.
% Mark as applicable:
% - actual pit stops;
% - tyre compounds/stint boundaries;
% - model decision point;
% - predicted gain/crossover;
% - actual post stop outcome;
% - excluded/invalid periods.
%
% SUGGESTED FIGURES:
% 4.2 unannotated historic race trace;
% 4.3 same trace with model annotations and explanatory callouts.
% Use real application output when available, not only mock-ups.

\subsection{Interaction and Usability}
% MEDIUM PRIORITY. Briefly explain race/session loading, driver selection and
% parameter options. Tie usability to rapid engineering interrogation rather
% than claiming a full formal UX study unless one was conducted.

\section{Software Verification}
\label{sec:software-verification}

% HIGH PRIORITY but concise. This belongs here/methodology Chapter 5, not as a
% dominant research result.
% Explain the test layers planned/used:
% - unit tests for individual equations and sign conventions;
% - boundary tests for tyre age, zero degradation and zero offset;
% - integration tests from data state to scenario result;
% - regression tests for known parameter sets;
% - manual comparison with the Chapter 3 worked example.
% Report actual results in Chapter \ref{chap:results} or an appendix.

\section{Development Process and Key Iterations}
\label{sec:development-process}

% Retain the useful diary/sprint narrative from the earlier report, but make it
% analytical rather than chronological padding. For each significant iteration:
% - intended outcome;
% - problem discovered;
% - evidence/constraint;
% - decision made;
% - consequence for model/research scope.
% Include revised dates following the leave of absence.
% Candidates: data-source changes, Blazor->WPF, API reliability, race trace
% development and model iteration.
% Detailed sprint logs may go to an appendix; keep only consequential decisions
% in the chapter.

\section{Implementation Limitations}
\label{sec:implementation-limitations}

% Discuss software/data limitations separately from model limitations:
% - dependency on third-party historic data;
% - Windows/WPF platform restriction;
% - manual parameters;
% - incomplete data or API reliability;
% - no live operational integration.
% State consequence and mitigation. Do not repeat Chapter 6's engineering model
% limitations verbatim.

\section{Chapter Summary}
\label{sec:implementation-summary}

% One concise paragraph linking the implemented workflow to the verification and
% validation methodology in Chapter \ref{chap:validation-method}.
